\section{算力扩散：Capacity、Inference 与数字鸿沟}

模型组合最终仍要落在基础设施上。课堂把大规模投资称为 no-regret，但系统上只有在 capacity 被利用、inference 可负担、应用可进入、收益可反馈时，这个判断才成立。本节用 flywheel 分析“买算力”怎样成功，也解释为什么它可能放大而非缩小数字鸿沟。

\subsection{Utilization flywheel 与 token economics}

上一节解决“选什么模型”，本节回答“怎样让国家级 capacity 真的被使用”。读图时需要区分 installed capacity、available service、实际 utilization 和 application demand：只有调度、模型 onboarding、可靠性和用户采用共同改善，unit cost 才可能下降，并反过来吸引更多有效 workload；单纯增加 rack 不会自动启动飞轮。

\lecturefigure{11-infrastructure-flywheel.png}{Capacity 通过 utilization、unit cost 与 application demand 才形成正反馈。}{本地字幕 27:10--32:20、42:30--45:10；Groq/Aramco Digital 官方背景。}

\begin{knowledgebox}{读图：最低 token price 不是唯一目标}
低价格可以扩大 experimentation，但用户还需要 model availability、Arabic/vertical quality、latency、data protection、reliability 和 support。若 workload 不稳定或软件栈难用，capacity 会变成补贴；若应用需求增长而 power 与 network 没有余量，低价承诺又可能损害 SLO。
\end{knowledgebox}

\begin{warningbox}{课堂时点与后续公告要分开}
访谈讨论的是 2025 年 3 月可见的 Groq/Aramco Digital 推理合作；之后的 LEAP 2025 扩展公告可用于说明演化，但不能倒写成课堂已经交付的容量。讲义因此只把“快速 onboarding 多模型、降低 inference 门槛”作为机制结论。
\end{warningbox}

\subsection{Training 与 inference 是两种基础设施形状}

\term{training}（训练）通常以大规模同步作业更新参数，追求吞吐和 time-to-train；\term{inference}（推理）用已训练模型响应请求，受 latency、availability、batching、KV cache、地域和持续成本约束。二者共享 accelerator，却不共享完全相同的 topology、scheduler 或 business demand。

\lecturefigure{12-training-inference.png}{Training 更集中和同步，inference 更分布且受 SLO 约束。}{本地字幕 32:20--36:50；概念重绘。}

\begin{knowledgebox}{读图：不要固定一个永久比例}
课堂对 training/inference 比例的判断是时点观察。模型研发、fine-tuning、synthetic data、test-time compute、agent loop 与用户增长都会改变 workload mix。Capacity plan 应用 scenario range 和 modular pool，而不是把当前比例外推到设施寿命末期。
\end{knowledgebox}

\begin{importantbox}{数字鸿沟的资源版与能力版}
资源版鸿沟是 access price、network 和 accelerator availability；能力版鸿沟是数据、人才、workflow redesign、evaluation 和 procurement。只降低 token price 可能改善前者，却不自动改善后者。公共计划需要把 adoption support 和 outcome measurement 纳入基础设施预算。
\end{importantbox}

\subsection{本章小结}

算力投资只有进入 utilization flywheel 才是 productive asset。Training 与 inference 的资源形状不同，数字普惠也不只是低价 token；software、skills、data 和 organizational adoption 同样是基础设施的一部分。

